%---------------------------------------------------------------------------------
%                西南交通大学研究生学位论文：英文摘要
%---------------------------------------------------------------------------------

% 英文摘要输入区
\eabstract
{
In the context of digital government construction, vast amounts of structured government affairs data have been continuously accumulated. 
This thesis investigates LLM-driven data query and analysis methods to address the pervasive challenges of data accessibility, analytical complexity, and collaborative silos faced by non-technical personnel in government scenarios.
Considering the complex mapping between government terminology and physical database fields, the implicit statistical calibers and computational logic, and the unclear boundary between query and analysis responsibilities, this thesis aims to build an intelligent data querying and analysis technology pipeline for the government affairs domain, guided by the overarching goal of ``reliable query, standardized analysis, and systematic implementation''.

Toward this goal, this thesis proposes a government affairs data query generation method based on logic enhancement and heterogeneous collaboration. By constructing a domain knowledge base, a dynamic example library, and a Logic-Augmented Schema (LA-Schema), business terminology mappings, logical computation formulas, and value domain constraints are explicitly injected into the query generation process. Combined with a logical mapping generator, a progressive divide-and-conquer generator, execution feedback-driven query repair, and a fine-tuned candidate discrimination model, the method achieves reliable query generation for complex government business semantics. Furthermore, this thesis proposes a government affairs data analysis method based on analysis planning and query augmentation. An Analysis Plan Schema (AP-Schema) is introduced to transform open-ended analysis requirements into structured plans. Through a query-augmented task decomposition mechanism, the reliable data retrieval capability is reused, and standardized analysis results are generated via lightweight analysis operator chains and visualization rules. Built upon these methods, this thesis designs and implements an intelligent data querying system for government scenarios, organizing the query and analysis methods into a unified, interactive, and traceable service.

Experimental results demonstrate that the proposed method achieves an execution accuracy of 91.47\% on the self-constructed government affairs dataset GovSQL, outperforming the strongest baseline XiYan-SQL by 7.64 percentage points. On government analysis tasks, the pipeline success rate, task completion rate, and result accuracy reach 92.9\%, 85.7\%, and 95.8\%, respectively, significantly outperforming mainstream analysis paradigms such as code generation and tool-iterative agents, with stable generation of charts, tables, and textual insights. The implemented system supports natural language data querying, analysis result generation, and process traceability, while maintaining strong competitiveness on the BIRD and Spider public datasets.

This research establishes a complete technology pipeline for the government affairs domain, spanning from natural language requirement understanding, reliable data retrieval, to standardized analysis and system implementation. It provides a methodological foundation and system support for non-technical personnel to conduct government data querying and analysis, and also offers a reference for the research and deployment of large language models in vertical domain data intelligence applications.
}
{Government Affairs Data, Large Language Model, Text-to-SQL, Intelligent Data Querying, Analysis Planning, Data Analysis}	% 英文关键词输入区，依据2024年新规范，关键词之间用半角逗号","分开，数量3~8个